% Generated draft: requires mathematical and source review.
\documentclass[11pt]{article}
\usepackage[margin=25mm]{geometry}
\usepackage{fontspec}
\setmainfont{Latin Modern Roman}
\usepackage{amsmath,amssymb,mathtools,unicode-math}
\setmathfont{Latin Modern Math}
\usepackage{enumitem,needspace,xurl,hyperref}
\setlength{\parindent}{0pt}
\setlength{\parskip}{5pt}
\setlength{\emergencystretch}{3em}
\newcommand{\N}{\mathbb{N}}
\newcommand{\Z}{\mathbb{Z}}
\newcommand{\Q}{\mathbb{Q}}
\newcommand{\R}{\mathbb{R}}
\newcommand{\C}{\mathbb{C}}
\newcommand{\F}{\mathbb{F}}
\newcommand{\A}{\mathbb{A}}
\newcommand{\PP}{\mathbb{P}}
\newcommand{\E}{\mathbb{E}}
\newcommand{\eps}{\varepsilon}
\newcommand{\dd}{\,\mathrm{d}}
\begin{document}
\section*{UnsolvedMath ID: 1100505}
\textbf{Source problem code:} catalogue entry
\subsection{Definitions and mathematical statement}
A Kleinian group in this question is a discrete subgroup $G\leq\operatorname{Isom}(\mathbb H^n)$; it is assumed finitely generated, but the statement imposes neither torsion-freeness nor geometric finiteness. The symbol $\operatorname{vcd}(G)$ conventionally denotes the cohomological dimension of a torsion-free finite-index subgroup, with the coefficient convention to be checked against the source. The record does not define $\delta(G)$. Under the usual critical-exponent convention, for $o\in\mathbb H^n$ it is
\[
 \delta(G)=\inf\left\{s\geq0:\sum_{g\in G}e^{-s d(o,go)}<\infty\right\}.
\]
This convention is explanatory, since the intended normalization of $\delta$ must be verified.

\textbf{Literal recorded question (Misha Kapovich).} For every finitely generated Kleinian group $G\leq\operatorname{Isom}(\mathbb H^n)$, is it true that
\[
 \delta(G)\leq\operatorname{vcd}(G),
 \qquad
 \delta(G)=\operatorname{vcd}(G)
 \Longleftrightarrow
 \exists\,\mathbb H^k\subseteq\mathbb H^n\text{ totally geodesic and $G$-invariant with }\mathbb H^k/G\text{ compact}?
\]
The source does not specify the permitted values of $k$ or exclude elementary groups. If torsion is present, the quotient is understood as the quotient of the action, possibly an orbifold.

\textbf{Incompatible accompanying formulation.} The supplied research summary instead reports
\[
 \delta(\Gamma)\geq\operatorname{vhd}_{\R}(\Gamma,\Pi)-1,
 \qquad
 \delta(\Gamma)\geq\operatorname{vcd}_{\R}(\Gamma,\Pi)-1
 \quad\text{when $(\Gamma,\Pi)$ has finite type}.
\]
It attributes these assertions for virtually torsion-free Kleinian groups to [Kap09], Theorem 1.1 and Corollary 1.2, and includes geometrically finite groups in the finite-type case. Here the summary uses relative virtual homological and cohomological dimensions over $\R$, but does not specify the family $\Pi$ or the precise finite-type convention. It further attributes to [Kap09], Theorem 1.3, the geometrically finite equality case: equality, described there as $\dim_{\mathrm{top}}\Lambda=\dim_H\Lambda=d$, forces the limit set $\Lambda$ to be a round $d$-sphere and $\Gamma$ to be a finite-volume lattice on a totally geodesic $\mathbb H^{d+1}$. It records the converse as $\delta=d=\operatorname{vcd}(\Gamma,\Pi)-1$ and says that equality implying geometric finiteness without assuming geometric finiteness remains open.

These reported statements have a different inequality direction, a dimension shift, relative dimensions, and a finite-volume conclusion where the literal question uses compactness. They therefore do not supply an unambiguous resolution or replacement of the displayed literal question. The normalization and intended formulation require source review.
\subsection{Short English statement}
Kapovich's recorded question compares an unspecified invariant $\delta(G)$ with virtual cohomological dimension and asks whether equality characterizes a compact action on an invariant totally geodesic hyperbolic subspace. The accompanying progress summary uses a different relative inequality and a finite-volume equality case, so the intended statement cannot be identified without reviewing the source.
\subsection{Sources}
\begin{itemize}
\item Local UnsolvedMath registry, permanent ID 1100505, source code catalogue entry, title ``Questions in Geometric Group Theory --- Q 5.5''; attributed in the statement to Misha Kapovich. Supplied status: partially solved; research classification: PARTIAL-PROGRESS; updated 2026-08-06.
\item Mladen Bestvina, \emph{Questions in Geometric Group Theory}, source list dated 2004, Question 5.5, PDF page 13; supplied URL \url{https://www.math.utah.edu/~bestvina/eprints/questions-updated.pdf}. The record reports manual transcription and access on 2026-07-29. Its background says that the author PDF records partial progress or special cases and that current status requires release review; rights note: NEEDS\_REVIEW, public author-hosted PDF.
\item The supplied research summary cites [Kap09], Theorem 1.1, Corollary 1.2, and Theorem 1.3, without a complete bibliographic entry. It also mentions unspecified conjecture citations from 2020--2022. These references and the summary's claimed results have not been independently verified here.
\end{itemize}
\end{document}
